\section{自测：把知识变成可以独立完成的判断}
\learning{先不看答案，写出依据或计算式；只记住最终数值不算完成。}

\subsection*{题目}
\begin{enumerate}[label=\textbf{\arabic*.},itemsep=0.65em]
\item 对 \ce{[CoCl(NH3)5]Cl2}，写出内界、外界、Co 的氧化态、$d$ 电子数和配位数。每摩尔该盐在未发生配体取代时，可立即提供多少摩尔外界氯离子？
\item 对六齿配位的 \ce{[Co(edta)]^-}，区分配体数、齿数与配位数。若只知道某配合物是四配位，是否足以断定它为正四面体？
\item 判断下列两组分别属于哪类异构：
\par \ce{[Co(NO2)(NH3)5]^2+} 与 \ce{[Co(ONO)(NH3)5]^2+}；
\par \ce{[CoBr(NH3)5]SO4} 与 \ce{[Co(SO4)(NH3)5]Br}。
\item 在忽略螯环构象、且双齿配体 AA 对称的理想模型中，八面体 $[M(AA)_2X_2]$ 有几种几何异构体、几种单分子立体异构体？给某一对对映体加入它们的等量混合物，是否增加一种分子异构体？
\item 某样品只含两个对映体，$R:S=97:3$。求 $ee$。一个配合物标为 $\Delta$，能否据此断定其旋光度为正？
\item 分别判断理想 \ce{H2O}、\ce{NH3}、\ce{BF3}、\ce{CH4}、\ce{SF6} 的点群。某固定构型没有镜面和反心，但有 $S_4$，它是否具有手性？
\item 在 $C_{2v}$ 中按 $E,C_2(z),\sigma(xz),\sigma(yz)$ 排列，求 $\Gamma_{xyz}$ 的特征标，并把 $(3,1,3,1)$ 约化。两个表示的维数为什么相同而变换性质不同？
\item 在 $C_{3v}$ 中，某个表示的特征标为 $(5,-1,1)$，表头为 $E,2C_3,3\sigma_v$。求它的不可约分解，并用维数检验。
\item 求 $C_{3v}$ 中 $E\otimes E$ 的分解。若跃迁被积函数的对称性包含 $A_1$，能否保证有强吸收？
\item 把水分子放在 $yz$ 平面，主轴沿 $z$。给定 $\Gamma_{3N}$ 的特征标为 $(9,-1,1,3)$，求振动表示、简正模式数和理想 IR 基频数。
\item \ce{NH3} 属 $C_{3v}$，$\Gamma_{3N}=(12,0,2)$。已知平移为 $A_1+E$、转动为 $A_2+E$，求振动表示、简正模式数和不同基频数；判断 IR、Raman 活性。
\item 为什么理想 $fac$-$ML_3(CO)_3$ 的三个 C--O 伸缩只有两个 IR 基频，而 $mer$ 型通常为三个？只分析这三根键时，是否需要再减去整体平移和转动？
\item 一个有反心的分子出现了 $g$ 对称性的振动，能否属于 IR 基频？所有 $u$ 振动是否必为 IR 活性？
\item 在 \ce{CH4} 的价层成键模型中，四个 H 的 $1s$ 轨道组成 $A_1+T_2$。它们分别与哪些 C 价轨道匹配？总共形成多少个价层 MO，多少个被占据？
\end{enumerate}

\clearpage
\subsection*{答案与回查}
\begin{enumerate}[label=\textbf{\arabic*.},itemsep=0.65em]
\item 内界为 \ce{[CoCl(NH3)5]^2+}，外界为两个 \ce{Cl^-}。$x-1=+2$，所以 Co 为 $+3$，$d^6$，配位数为 $6$。每摩尔盐可立即提供 $2\ \mathrm{mol}$ 外界 \ce{Cl^-}。\textcolor{Muted}{回查第1节。}
\item 配体数为 $1$，edta 的实际齿数为 $6$，配位数为 $6$。四配位还可能是平面四方形，单靠配位数不能判断。\textcolor{Muted}{回查第1--2节。}
\item 前一对是键合异构，改变 N/O 配位原子；后一对是电离异构，内外界发生交换。\textcolor{Muted}{回查第3节。}
\item 有 $cis$、$trans$ 两种几何异构体；$cis$ 有一对对映体，$trans$ 在这个模型中非手性，所以共有 $3$ 种单分子立体异构体。等量混合物不增加分子异构体数。\textcolor{Muted}{回查第3节。}
\item $ee=(97-3)/(97+3)\times100\%=94\%$，$R$ 过量。$\Delta$ 是构型标签，不能推出旋光正负。\textcolor{Muted}{回查第3--4节。}
\item 依次为 $C_{2v},C_{3v},D_{3h},T_d,O_h$。有 $S_4$ 的固定构型不具有手性，是否存在非正旋转操作才是完整检查。\textcolor{Muted}{回查第5节。}
\item $\Gamma_{xyz}=(3,-1,1,1)=A_1+B_1+B_2$。另一个表示中
\[
 a_{A_1}=\tfrac14(3+1+3+1)=2,\qquad
 a_{B_1}=\tfrac14(3-1+3-1)=1,
\]
其余系数为零，所以 $(3,1,3,1)=2A_1+B_1$。两者都由三个基组成，但基在各操作下的变化不同。\textcolor{Muted}{回查第6--7节。}
\item $a_{A_1}=(5-2+3)/6=1$，$a_{A_2}=(5-2-3)/6=0$，$a_E=(10+2)/6=2$。因此为 $A_1+2E$，维数 $1+2\times2=5$。\textcolor{Muted}{回查第7节。}
\item $E\otimes E=A_1+A_2+E$。包含 $A_1$ 只表示对称性不强制积分为零，强度还取决于具体矩阵元等因素。\textcolor{Muted}{回查第7节。}
\item 约化得 $\Gamma_{3N}=3A_1+A_2+2B_1+3B_2$；减去平移 $A_1+B_1+B_2$ 和转动 $A_2+B_1+B_2$，得
\[
 \Gamma_{\mathrm{vib}}=2A_1+B_2.
\]
共 $3$ 个模式、$3$ 个不同基频，均有 IR 活性。\textcolor{Muted}{回查第8节。}
\item $\Gamma_{3N}=3A_1+A_2+4E$，故 $\Gamma_{\mathrm{vib}}=2A_1+2E$。共有 $2+2\times2=6$ 个模式、$4$ 个不同基频；$A_1$ 与 $E$ 都出现在坐标及二次函数中，因此四个基频均可有 IR 和 Raman 活性。\textcolor{Muted}{回查第7--8节。}
\item $fac$ 型伸缩表示为 $A_1+E$，$E$ 二重简并，所以为两个 IR 基频；$mer$ 型为两个全对称伸缩加一个一维反对称伸缩，三者均可有 IR 活性。伸长量已是内部坐标，不再减平移、转动。\textcolor{Muted}{回查第9节。}
\item $g$ 模式不能是电偶极 IR 基频，因为 $x,y,z$ 都是 $u$。但 $u$ 只是必要条件：还必须与至少一个坐标分量同型；否则可以是静默模式。\textcolor{Muted}{回查第8节。}
\item $A_1$ 与 C $2s$ 配对，$T_2$ 与 C $(2p_x,2p_y,2p_z)$ 配对。四个 H 轨道加四个 C 价轨道，共形成八个价层 MO：四个成键、四个反键；八个价电子占据四个成键 MO。\textcolor{Muted}{回查第10节。}
\end{enumerate}

\clearpage
\section{一页速查：从题目出发选方法}
\normalsize
\noindent
\begin{tabularx}{\linewidth}{@{}p{30mm} X@{}}
\toprule
题目给了什么 & 下一步做什么 \\
\midrule
配合物化学式 & 用“金属氧化态 $+$ 配体电荷总和 $=$ 络离子电荷”；再数实际配位原子，不能用配体个数代替配位数。过渡金属常规计数 $d^n$：$n=\text{族数}-\text{氧化态}$。\\
配位数与配体 & 四配位重点分正四面体、平面四方形；五配位分三角双锥、四方锥；六配位重点看八面体，同时记住三棱柱可能存在。\\
一组同分异构体 & 先查连接关系：变了是结构异构，没变再查空间排列。八面体 $MA_4B_2$ 判 $cis/trans$；$MA_3B_3$ 判 $fac/mer$。\\
手性与旋光信息 & 固定构型无任何 $S_n$ 才有手性；$\sigma=S_1$、$i=S_2$。$\Delta/\Lambda$ 不决定 $(+)/(−)$。对映体混合物不另算一种分子异构体。\\
对映体比例 & $ee=|n_R-n_S|/(n_R+n_S)\times100\%$；多数体摩尔分数 $=(1+ee/100\%)/2$。$dr$ 比较非对映体，不能替代 $ee$。\\
需要判点群 & 先查线形与高对称结构；再找最高次 $C_n$；检查垂直 $C_2$ 轴，再查水平/竖直镜面与反轴。把配体骨架和取代基一起算入。\\
一行可约特征标 & 约化：$a_i=h^{-1}\sum_k n_k\chi_{\mathrm{red}}(R_k)\chi_i(R_k)^*$。系数应为非负整数；用 $\sum_i a_il_i=\chi_{\mathrm{red}}(E)$ 检查维数。\\
全部振动模式 & $\chi_{3N}(R)=N_{\mathrm{fixed}}(R)\chi_{xyz}(R)$。约化后减平移与转动；非线形总数 $3N-6$，线形 $3N-5$。\\
特定一组键伸缩 & 用各键的伸长量作基：每个操作下数保持原位的目标键。直接约化伸缩表示，不乘 $\chi_{xyz}$，也不减整体平移、转动。\\
IR / Raman 谱带 & IR 查 $x,y,z$；Raman 查 $x^2,y^2,z^2,xy,xz,yz$。按不可约表示的出现次数数不同基频，简并维数计入模式数而不重复计频率。\\
有反演中心 & $g$ 模式可有 Raman 活性，$u$ 模式可有 IR 活性；同一模式不能两者兼有。并非所有 $g/u$ 模式都一定活跃。\\
跃迁或轨道耦合 & 逐列乘特征标再约化。$\Gamma_f^*\otimes\Gamma_O\otimes\Gamma_i$ 不含全对称，矩阵元必为零；含全对称只是对称允许。轨道还需能量接近、有效重叠。\\
\bottomrule
\end{tabularx}
\par\vspace{0.8em}
\noindent\textcolor{Muted}{\textbf{最后一次检查}\quad 电荷是否守恒？几何模型是否说清？约化维数是否守恒？数的是模式、频率，还是实际可分辨的峰？}
\normalsize
